\newpage
\begin{framed}
\noindent\textbf{Box 1. Linking variance across discrete and continuous time models} 

During the early stages of an outbreak, secondary-cases counts are typically modeled using a negative-binomial distribution, with mean $\mu$ and shape, or “dispersion,” parameter $k$,. 
This distribution corresponds mechanistically to a gamma-distributed  “intrinsic infectiousness”, with the same mean and shape parameters, and a Poisson distribution of realized cases based on this infectiousness \cite{lloydsmith2005superspreading}.

This is illustrated in \fref{ls} for the case where intrinsic infectiousness arises from random recovery times under simple SIR assumptions, and is thus exponentially distributed. Expected infectiousness scales exactly with \Ro, so that the smooth curves in panel b, are all exactly the same, but the proportional amount of randomness added by the Poisson step increases when \Ro\ is small.

We note a couple of points. First, the squared coefficient of variation (SCV) of the underlying Gamma distribution is $\kappa=1/k$, so its variance is $\sigmaAct^2 = \kappa \mu^2$. The resulting variance in secondary-case counts is $\sigma^2 = \kappa\mu^2 + \mu$, with the last term arising from the Poisson step of the negative-binomial construction.
In the case of \fref{ls} -- exponentially distributed infectiousness (corresponding to a geometric distribution of secondary cases) -- $\kappa = k = 1$.
Second, the relationship between $\sigmaAct^2$ and $\sigma^2$ holds regardless of the form of the expected infectiousness distribution, due to the Law of Total Variation \cite{mixedPoissonDist}. JD: Check does this cite mention the Law?

We therefore suggest $\kappa$ as a generalized dispersion parameter. This has the advantage of less dependence on a specific assumption about distribution (and parameterization), and of the intuitive idea that variation increases when “dispersion” increases.


\end{framed}
